% Noether Paper 20 controlled-generic zh-Hant producer translation.
% Current-German interval SHA-256:
% CBC9E9CF34E6475F4256C935A58378FCDBF85A09ACC0E592FC64F3FCFDF8744D.
% Inherited Chinese content witness SHA-256:
% B7DA9DBB83BC2B9793263987F14D7E67C91324EE83FEA26FBCDC82979FE5F97C.
% Controlled-generic script transport only; every independent check pending.
% Controlled generic Traditional script only; not zh-Hant-TW/HK/MO prose.
% Hans wording remains the lexical base; independent Hant checking is pending.
% Compilation is mechanical production only; no source, semantic, formula-content,
% terminology, visual, regional, publication, archive, or certification check is claimed.
\documentclass[11pt]{article}
\usepackage[a4paper,margin=2.4cm]{geometry}
\usepackage{fontspec}
\usepackage{xeCJK}
\usepackage{amsmath,amssymb,mathtools,array}
\setmainfont{Noto Serif}
\setCJKmainfont{Microsoft JhengHei}
\setCJKsansfont{Microsoft JhengHei}
\XeTeXlinebreaklocale "zh"
\XeTeXlinebreakskip=0pt plus 1pt
\setlength{\parindent}{2em}
\setlength{\parskip}{0.35em}
\setlength{\emergencystretch}{8em}
\allowdisplaybreaks
\raggedbottom
\providecommand{\glqq}{“}
\providecommand{\grqq}{”}
\providecommand{\frK}{\mathfrak K}
\providecommand{\srcfn}[2]{%
  \begingroup
  \renewcommand{\thefootnote}{#1}%
  \footnote{#2}%
  \addtocounter{footnote}{-1}%
  \endgroup}
\providecommand{\srcfnmark}[1]{%
  \begingroup
  \renewcommand{\thefootnote}{#1}%
  \footnotemark%
  \addtocounter{footnote}{-1}%
  \endgroup}
\providecommand{\srcfntext}[2]{%
  \begingroup
  \renewcommand{\thefootnote}{#1}%
  \footnotetext{#2}%
  \endgroup}
\providecommand{\srcnumdisplay}[3][3.4em]{%
  \par\smallskip\noindent
  \begin{minipage}[t]{#1}\vspace*{0pt}#2\end{minipage}%
  \begin{minipage}[t]{\dimexpr\linewidth-#1\relax}\vspace*{0pt}\centering
  \(\displaystyle #3\)
  \end{minipage}\par\smallskip}
\begin{document}
\setcounter{footnote}{0}

\section*{20. 絕對不可約性的一個代數判據}
\begin{center}
作者\\[0.25em]
哥廷根的埃米·诺特。\\[0.6em]
Math. Ann. 85（1922），第26--33页
\end{center}

设一个多项式的系数属于任意一个抽象定义的域。如果将系数域扩张到
一个代数闭域后，该多项式仍然不可约，就称它为\emph{绝对不可约的}
\srcfnmark{1)}。\srcfntext{1)}{E. Steinitz 在他的《域的代数理论》
（J. f. M. 137 (1910)，第167页）中证明了，每个域都能以本质上唯一的
方式扩张为代数闭域，即扩张为一个使每个一元多项式都能分解为一次因子的
域；由此他也给出了代数基本定理的有理对应形式。}\addtocounter{footnote}{1}
下文我将说明，与相对于一个给定域的不可约性不同，绝对不可约性可以
\emph{用代数方式}刻画；更确切地说，有如下定理：

\emph{对于每个具有 \(n\ge2\) 个变量和未定系数的多项式，都可以给它
对应一个关于这些系数和另外一些不定元的整有理、整系数函数，即可约性
形式；对于每个同次数的特殊多项式，绝对不可约的充要条件就是这个可约性
形式不恒等于零。} 换言之：\emph{对于每一组特殊的系数值，如果可约性
形式关于各不定元恒等于零，并且这组值不导致降次，那么相应的特殊多项式
就是可约的；反之亦然。}

如果不用多项式，而考察多一个变量的齐次形式，那么次数条件就由一个
更简单的条件取代，即系数系不恒等于零。

作为这个定理的直接推论，还可得到 A. Ostrowski 最先提出的结论
\footnote{《论代数量的算术理论》。Gött. Nachr. 1919，第279页
（引理见第296页）。-- 据我所知，对于由通常数构成的域，Ostrowski
也得到了上面的主定理；他将于近期发表有关研究。在我的证明安排中，
主定理对任意抽象定义的域都成立，这是因为对不定元构成的可约性形式
具有\emph{与具体选作基础的域无关的数系数}。}：

\emph{每个以代数数为系数的绝对不可约多项式，模包含其系数域的任意
有限代数域 \(\Omega\) 的每个素理想仍然不可约，至多只有 \(\Omega\)
中的有限多个素理想例外。特别地，\(\Omega\) 也可以是有理数域。}

关于相对整函数理论的进一步应用\footnote{D. Hilbert，《数学问题》。
Gött. Nachr. 1900，第253页，问题14。}，我拟在别处讨论。

1. 首先证明：\emph{每个具有 \(n\ge2\) 个变量和未定系数的多项式都
是绝对不可约的。} 令
\srcnumdisplay{(1)}{%
\begin{aligned}
F(x)&=\sum A_{i_1\ldots i_n}x_1^{i_1}\cdots x_n^{i_n} &&(i_1+\cdots+i_n\le l),\\
G(x)&=\sum B_{j_1\ldots j_n}x_1^{j_1}\cdots x_n^{j_n} &&(j_1+\cdots+j_n\le m),\\
H(x)&=F(x)G(x)=\sum C_{k_1\ldots k_n}(A,B)x_1^{k_1}\cdots x_n^{k_n} &&(k_1+\cdots+k_n\le l+m),
\end{aligned}
}
其中 (A) 和 (B) 表示不定元，而 (C(A,B)) 是 (A) 和 (B) 的
整系数双线性组合。因此，商
\(D=C(A,B)/C_{0\ldots0}(A,B)\) 是各商
\(A/A_{0\ldots0}\) 和 \(B/B_{0\ldots0}\) 的有理函数；但是这些函数的
个数大于其宗量的个数，因为前者为
\(\binom{l+m+n}{n}-1\)，而两组宗量的个数分别为
\(\binom{l+n}{n}-1\) 和 \(\binom{m+n}{n}-1\)，并且有
\[
\binom{l+m+n}{n}+1>
\binom{l+n}{n}+\binom{m+n}{n}
\quad\text{当 } n\ge2,\ l>0,\ m>0\text{ 时}.
\]
因此，諸 (D(A,B)) 之間至少存在\emph{一個}代數關係，從而諸
\(C(A,B)\) 之間也存在一個代數關係：
\srcnumdisplay{(2)}{%
\Phi(Z)\ne0\quad[\text{關於 }Z_{k_1\ldots k_n}\text{ 恆等}],\ 
\Phi(C(A,B))=0\quad[\text{關於 }A,B\text{ 恆等}],
}
其中 \(\Phi\) 表示一個整係數多項式。

因此，以\emph{不定元 (Z)} 為係數的多項式
\srcnumdisplay{(3)}{%
E(x)=\sum Z_{k_1\ldots k_n}x_1^{k_1}\cdots x_n^{k_n}
\qquad(k_1+\cdots+k_n\le l+m)
}
在 \(n\ge2\) 時必定\emph{絕對不可約}\srcfnmark{4)}。
\srcfntext{4)}{因為按照第2節的說法，不定元 (Z) 不是
\(\mathfrak P\) 的零點。}\addtocounter{footnote}{1}

2. 在代換 \(Z=C(A,B)\) 後關於 (A,B) 恆等於零的所有這些多項式
\(\Phi(Z)\) 構成一個\emph{多項式理想}\srcfnmark{5)}，
\srcfntext{5)}{這些總體通常稱為模；但真正刻畫它們的性質是：它們在
包含 \(\Phi_1\) 和 \(\Phi_2\) 時也包含 \(\Phi_1-\Phi_2\)，在包含
\(\Phi\) 時也包含 \(\Psi\Phi\)，其中 \(\Psi\) 是 (Z) 的任意整係數
多項式；這實際上是理想性質。}\addtocounter{footnote}{1}
更確切地說，構成一個\emph{素理想 \(\mathfrak P\)}；因為如果藉助這個
代換，一個乘積 \(\Phi_1\Phi_2\) 恆等於零，那麼至少有一個因子恆等於
零。由於(2)關於 (A,B) 恆等成立，所以對於每一組特殊值 (A,B)，
(2)仍然成立；這裡 (A,B)，因而 \(\Gamma=C(A,B)\)，可以是任意一個
域中的量。\emph{所以，在某個擴域中可約的每個多項式的係數
\(\Gamma\) 都滿足條件 \(\Phi(\Gamma)=0\)，即為 \(\mathfrak P\) 的
零點，}也就是 \(\mathfrak P\) 的有限多個基多項式的公共零點；現在還須
\emph{證明逆命題}。

為此要注意，如果通過引入一個新變量 (y)，把(1)中的多項式變為
次數分別為 (l,m,l+m) 的齊次形式 (F(x,y),G(x,y),H(x,y))，那麼
 (A,B,C) 之間的一切關係仍然保持不變。因此，如果某些係數
\(\Gamma\) 使例如 \(F(x,y)\)\footnote{把
\(F,G,\ldots,f,g,\ldots\) 中的不定元換成特殊值系而得到的多項式
和形式，下文一律以加上橫線來表示：
\(\bar F,\bar G,\ldots,\bar f,\bar g,\ldots\)。} 成為 (y) 的一個
冪，那麼這些係數也成為 \(\mathfrak P\) 的零點；而從齊次形式轉回非齊次
多項式時，與此相應的是 \(\bar H(x)\) 的降次，並不對應於分解。下文將
證明，除這種降次外，逆命題總是成立；特別是，\emph{對於齊次形式，只要
各 \(\Gamma\) 不全為零，逆命題就毫無例外地成立；換言之，
\(\mathfrak P\) 的每個零點都能由參數表示 \(\Gamma=C(A,B)\) 給出，
其中 (A,B) 是某個擴域中的量}\srcfnmark{7)}。
\srcfntext{7)}{由素理想 \(\mathfrak P\) 定義一個不可約代數構形這一
性質，可以推出 \(\mathfrak P\) 的零點“在一般情形下”由參數表示給出。
但眾所周知，這並不能推出參數表示對於每一組特殊值都不失效；我最初忽略
了這一點，A. Ostrowski 早已向我指出。這裡給出的證明建立在更初等的
概念上。}\addtocounter{footnote}{1}

我將通過直接構造有限多個函數 \(\Phi(Z)\) 來完成這個證明。具體說，
藉助 Kronecker 代換
\[
 x_\lambda=\xi^{d^{\,n-\lambda}}\text{\srcfnmark{8)}}
\]
我把(1)中的多項式對應到\emph{一個}變量的多項式；指出其指數中會出現
空缺，再通過構造相應係數的範數，得到函數 \(\Phi(Z)\) 以及所求判據。

\srcfntext{8)}{Festschrift，第11頁。}
\addtocounter{footnote}{1}
3. 令：
\srcnumdisplay{(4)}{%
\begin{aligned}
d&=l+m+1; &\qquad i&=i_1d^{n-1}+i_2d^{n-2}+\cdots+i_n;\\
j&=j_1d^{n-1}+\cdots+j_n; &\qquad k&=k_1d^{n-1}+\cdots+k_n.
\end{aligned}
}
藉助 Kronecker 代換，(1) 中的多項式 $F,G,H$ 分別對應於多項式
\srcnumdisplay{(5)}{%
\begin{aligned}
 f(\xi)&=\sum A_{i_1\ldots i_n}\xi^i &&(i_1+\cdots+i_n\le l),\\
 g(\xi)&=\sum B_{j_1\ldots j_n}\xi^j &&(j_1+\cdots+j_n\le m),\\
 h(\xi)&=\sum C_{k_1\ldots k_n}(A,B)\xi^k &&(k_1+\cdots+k_n\le l+m).
\end{aligned}
}
眾所周知，這個對應是\emph{一一對應}，因為對於 (5) 中出現的所有
\emph{（誘導）}指數，由 $k=0$ 也可推出
$k_1=0,\ldots,k_n=0$；因而由 $i+j=i'+j'$ 還可推出
$i_1+j_1=i'_1+j'_1;\ldots;i_n+j_n=i'_n+j'_n$。所以，只有當相應的
$x_1^{i_1}\cdots x_n^{i_n}x_1^{j_1}\cdots x_n^{j_n}$ 也產生相同的
指數時，不同的乘積 $\xi^i\xi^j$ 才可能產生同一個指數 $\xi^k$。
於是得到：
\srcnumdisplay{(6)}{%
h(\xi)=f(\xi)g(\xi),
}
而且，如果關係 (6) 成立，並且 $f(\xi)$ 與 $g(\xi)$ 中不出現不同於
誘導指數的指數，那麼反過來可得
\srcnumdisplay{(7)}{%
H(x)=F(x)G(x).
}
這裡，\emph{$f(\xi)$ 與 $g(\xi)$ 的誘導指數中確實存在空缺}。
事實上，$f(\xi)$ 的最高指數為 $ld^{n-1}$，次高指數為
$(l-1)d^{n-1}+d^{n-2}$；二者之差
$d^{n-2}(d-1)>1$，因為 $d=l+m+1\ge3$ 且 $n\ge2$；對於
$g(\xi)$ 也是如此。

對不定元 $A,B$ 成立的關係 (6) 與 (7)，對於 $A,B$ 的每個特定
值組仍然成立。為了使次數保持不變，我藉助變元 $\eta$ 與 $y$
將 (6) 和 (7) 齊次化。\emph{因此，齊次形式 $H(x,y)$ 在某個
代數擴張域中分解，當且僅當相應的二元形式 $h(\xi,\eta)$ 在該域中
能夠分解為兩個因子，並且這些因子中不出現不同於誘導指數的指數。}

4. 為了實施第 2 節末尾所述的範數構造，令
\[
 a_1,b_1;\ldots; a_p,b_p;\quad a_{p+1},b_{p+1};\ldots; a_{p+q},b_{p+q}
\]
表示新的不定元。令
\srcnumdisplay{(8)}{%
\begin{aligned}
 \varphi(\xi,\eta)&=(a_1\xi+b_1\eta)\cdots(a_p\xi+b_p\eta)
     =r_p\xi^p+r_{p-1}\xi^{p-1}\eta+\cdots+r_0\eta^p,\\
 \psi(\xi,\eta)&=(a_{p+1}\xi+b_{p+1}\eta)\cdots(a_{p+q}\xi+b_{p+q}\eta)
     =s_q\xi^q+s_{q-1}\xi^{q-1}\eta+\cdots+s_0\eta^q,\\
 \chi(\xi,\eta)&=\varphi(\xi,\eta)\psi(\xi,\eta)
     =t_{p+q}\xi^{p+q}+t_{p+q-1}\xi^{p+q-1}\eta+\cdots+t_0\eta^{p+q};
\end{aligned}
}
這裡 $r,s,t$ 分別表示齊次化的基本對稱函數；也就是說，它們是關於
$a,b$ 各列分別為線性的對稱函數，並且每個在各列中齊次且同維數的
整係數對稱函數，都能由它們整式地且以整數係數構成，而且這種構成
方式是唯一的。

現在引入另外的不定元 $u_{\mu\nu}$，並構造線性形式
\srcnumdisplay{(9)}{%
\zeta(u,a,b)=\sum u_{\mu\nu}r_\mu s_\nu,
}
其中求和遍及預先給定的指標 $\mu$ 與 $\nu$。將 $\zeta(u)$ 與所有
通過置換序列 $a,b$ 而產生、且與 $\zeta(u)$ 及彼此均不相同的共軛
量相乘——也就是說，把 $\zeta(u)$ 看作由 $r_\mu(a,b)$ 與
$s_\nu(a,b)$ 所生成的域中的函數時，取它關於由 $t(a,b)$ 所生成
的域的範數 $N(\zeta(u))$——所得結果是所有序列 $a,b$ 的一個
整係數對稱函數，它在每個序列中齊次且同維數。因此，根據上面的
結論，它是 $t(a,b)$ 與 $u_{\mu\nu}$ 的一個唯一確定的整係數
整函數：
\srcnumdisplay{(10)}{%
N(\zeta(u,a,b))=T(t(a,b),u)\qquad[\text{關於 }a,b,u\text{ 恆等}].
}
同樣的論證表明，還可得到
\[
N(z-\zeta(u))=z^\delta+T_{\delta-1}(t,u)z^{\delta-1}+\cdots+T(t,u)
\qquad[\text{关于 }a,b,u,z\text{ 恒等}],
\]
其中所有 $T$ 都表示 $t$ 與 $u$ 的整係數整函數。這個恆等式的
兩邊在 $z=\zeta(u)$ 時都消失；而且，由於基本對稱函數
$r(a,b),s(a,b)$ 在代數上獨立，當依照 (8) 將 $t(a,b)$ 置為
$r$ 與 $s$ 的一個整係數雙線性組合 $w(r,s)$，並將 $\zeta(u)$
視為 $r$ 與 $s$ 的函數時，這種消失關於 $r,s$ 恆等成立。因此
得到
\srcnumdisplay{(11)}{%
\zeta(u)^\delta+T_{\delta-1}(w(r,s),u)\zeta(u)^{\delta-1}+\cdots+T(w(r,s),u)=0
\quad[\text{關於 }r,s,u\text{ 恆等}].\,\text{\textsuperscript{9)}}
}
\srcfntext{9)}{當 (9) 中的求和遍及所有指標時，(11) 表示 Kronecker
對 Gauss 定理的推廣，而且本質上採用了 Kronecker 的證明編排
（Zur Theorie der Formen höherer Stufen, Berliner Ber. 1883,
Werke, 2, S.~417）。}
\addtocounter{footnote}{1}

恆等式 (10) 與 (11) 對於 $a,b$ 的所有特定值組 $\alpha,\beta$，
或者相應地對於 $r,s$ 的所有特定值組 $\varrho,\sigma$，仍然成立。
特別是，如果選取 $\varrho,\sigma$，使所有乘積
$\varrho_\mu\sigma_\nu$ 都消失——這裡 $\mu,\nu$ 表示 (9) 中出現的
指標——從而 $\zeta(u)$ 在代入這些值後消失，那麼令
$w(\varrho,\sigma)=\tau$，由 (11) 得：
\srcnumdisplay{(12)}{%
T(\tau,u)=0\qquad[\text{關於 }u\text{ 恆等}].
}

反過來，設 $\tau_0,\ldots,\tau_{p+q}$ 是不全為零的
\emph{值組，並且滿足 (12)。} 由於在某個代數擴張域中存在分解
\srcnumdisplay{(12)}{%
\bar\chi(\xi,\eta)=\tau_{p+q}\xi^{p+q}+\cdots+\tau_0\eta^{p+q}
=(\alpha_1\xi+\beta_1\eta)\cdots(\alpha_{p+q}\xi+\beta_{p+q}\eta),
}
其中任何一個因子中的 $\alpha$ 與 $\beta$ 都不同時為零
\srcfnmark{10)},\srcfntext{10)}{這個（有理的）“代數學基本定理”
是一個存在性定理，第 2 節末尾所說的逆命題無例外成立，正是以它
為依據。}\addtocounter{footnote}{1} 但某些 $\alpha$ 或 $\beta$
可以為零，於是 $\tau=t(\alpha,\beta)$。把這些值代入 (10)，
並利用 (12)，可知 $\zeta(u,\alpha,\beta)$ 或它的某一個共軛
因子關於 $u$ 恆等消失。因此，在適當重新編號 $\alpha,\beta$ 後，
若令 $r(\alpha,\beta)=\varrho$、$s(\alpha,\beta)=\sigma$，
這就意味着 $\bar\chi(\xi,\eta)$ \emph{至少容許一種分解：}
\srcnumdisplay{(13)}{%
\bar\chi(\xi,\eta)=\bar\varphi(\xi,\eta)\bar\psi(\xi,\eta)
=\sum \varrho_\kappa\xi^\kappa\eta^{p-\kappa}\cdot\sum \sigma_\lambda\xi^\lambda\eta^{q-\lambda},
}
\emph{使得 $\zeta(u)$ 中出現的所有乘積
$\varrho_\mu\sigma_\nu$ 都消失，但 $\varrho$ 或 $\sigma$ 並非
全部為零。}

5. 現在令 (8) 中的次數 $p,q$ 分別等於 $ld^{n-1}$ 和 $md^{n-1}$，即等於 (5) 中 $f(\xi)$ 和 $g(\xi)$ 的次數。把指標 $\mu,\nu$ 分拆為 $\mu^{(1)},\nu^{(1)},\mu^{(2)},\nu^{(2)}$；其中 $\mu^{(1)}$ 遍歷 $f(\xi)$ 中缺失的所有指數，$\nu^{(1)}$ 遍歷 $\psi(\xi,\eta)$ 中 $\xi$ 的所有指數；相應地，$\nu^{(2)}$ 遍歷 $g(\xi)$ 中缺失的所有指數，$\mu^{(2)}$ 遍歷 $\varphi(\xi,\eta)$ 中 $\xi$ 的所有指數。再令
\[
 e(\xi)=\sum Z_{k_1\ldots k_n}\xi^k\qquad(k_1+\cdots+k_n\le l+m)
\]
表示與多項式 (3) 相應的一元多項式；
\[
 S(Z,u)
\]
表示關於 $Z$ 和 $u$ 的整係數多項式；它由 $T(t,u)$ --- 即 (10) 中唯一確定的右端，其中 (9) 的求和遍及上述指標 $\mu^{(1)},\nu^{(1)},\mu^{(2)},\nu^{(2)}$ --- 通過把 $t$ 替換為係數 $Z$ 以及 $e(\xi)$ 的相應零點而得到。

現在把 $r,s$ 替換為 $f(\xi)$ 和 $g(\xi)$ 的係數 $A,B$ 以及相應零點，則在指定的求和下，經此代入後 $\zeta(u)$ \emph{消失}。此外，$t=w(r,s)$ 變為 $h(\xi)$ 的係數 $C(A,B)$ 及其相應零點；因此 $T(t,u)$ 變為 $S(C(A,B),u)$。由恆等式 (11) 便得到
\srcnumdisplay{(14)}{%
S(C(A,B),u)=0.\qquad[\text{關於 }A,B,u\text{ 恆等}].
}
由於 (14) 對每一個特化值系都保持成立，所以在一個\emph{擴張域中分解為兩個因子}的這種特化 $\bar H(x)$ 或一般的 $\bar H(x,y)$，其係數 $\Gamma=C(A,B)$ 滿足條件
\srcnumdisplay{(15)}{%
S(\Gamma,u)=0\qquad[\text{關於 }u\text{ 恆等}].
}

反過來，如果多項式 $\bar H(x)$，或形式 $\bar H(x,y)$ 的不全為零的係數 $\Gamma$ 滿足條件 (15)，那麼根據上面的論述，這等價於 $T(\tau,u)=0$，其中 $\tau$ 表示 $\bar h(\xi,\eta)$ 的全部係數。因而根據 (13)，在 $\Gamma$ 的一個擴張域中存在分解：
\[
 \bar h(\xi,\eta)=\bar f(\xi,\eta)\cdot\bar g(\xi,\eta),
\]
其中，由於指定的求和，$\bar f$ 和 $\bar g$ 中不出現任何不同於誘導指數的指數；於是存在關係
\[
 \bar H(x,y)=\bar F(x,y)\cdot\bar G(x,y).
\]
由此，只要在 $y=1$ 時沒有因子變為零次，也就是說，特別是在 $\Gamma$ 不會導致 $\bar H(x)$ 次數降低時，還可得到關係
\[
 \bar H(x)=\bar F(x)\cdot\bar G(x).
\]

\emph{因此，條件 (15) 是 $\bar H(x,y)$ 在一個擴張域中分解為兩個因子的充要條件；在排除次數降低時，它也是 $\bar H(x)$ 作這種分解的充要條件。}

由此還可推出，$S(Z,u)$ 不可能關於 $Z$ 和 $u$ 恆等消失，因為第 1 節已經證明，當 $n\ge2$ 時，以不定元為係數的多項式，因而也包括相應的多一個變量的齊次形式，都是絕對不可約的。因此，若用 $U$ 表示 $u$ 的冪積，則有：
\[
S(Z,u)=\sum \Phi_\lambda(Z)U_\lambda,
\]
其中 $\Phi_\lambda(Z)$ 根據 (14) 成為屬於 $\mathfrak P$ 的多項式。\emph{由此證明了，$\mathfrak P$ 除了由參數表示 $\Gamma=C(A,B)$ 給出的零點以外，沒有其他零點；這裡 $A$ 和 $B$ 屬於 $\Gamma$ 的一個代數擴張域。}這樣便證明了第 2 節末尾所陳述的逆命題。

6. \emph{現在可以直接表述絕對不可約性的條件。}為此，只需把 (3) 中 $E(x)$ 的次數以所有可能的方式分解為兩個正數之和 $l+m$，併為每個分解構造形式 $S(Z,u)$。這些形式的乘積
\srcnumdisplay{(16)}{%
R(Z,u)=\prod S(Z,u)
}
便構成 $E(x)$ 的\emph{可約性形式}；因為根據上面的論述，$\bar E(x)$ 或 $\bar E(x,y)$ 的每一次\emph{分解}，都對應於 (16) 的一個因子消失，反過來也成立。這樣便證明了\emph{引言中提出的主定理}，同時也說明了如何用有限多個步驟實際構造可約性形式。

7. 由\emph{可約性形式的存在性} (16)，立即可以推出引言中提到的\emph{奧斯特羅夫斯基定理}。設
\[
 \bar E(x)=\sum \Gamma_{k_1\ldots k_n}x_1^{k_1}\cdots x_n^{k_n}
 \qquad(k_1+\cdots+k_n\le v)
\]
是一個以代數整數 $\Gamma$ 為係數的絕對不可約多項式；$\frK$ 表示包含所有 $\Gamma$ 的一個有限代數數域，$\mathfrak p$ 表示 $\frK$ 中的一個素理想。

那麼，為 $\bar E(x)$ 的準確次數構造的可約性形式 $R(Z,u)$，在代入 $Z=\Gamma$ 後仍然\emph{不為零}；即
\[
 R(\Gamma,u)=\sum P_\lambda U_\lambda\ne0\qquad[\text{关于 }u\text{ 恒等}].
\]
因而，理想 $\mathfrak r=(P_1,P_2,\ldots)$ 只能被 $\frK$ 中\emph{有限多個}素理想整除。因此，如果 $\mathfrak p$ 不同於這有限多個素理想，那麼 $R(\Gamma,u)$ 模 $\mathfrak p$ 不會恆等消失；又因為模 $\mathfrak p$ 的剩餘類仍然構成一個域，所以可以推出，$\bar E(x)$ \emph{模 $\mathfrak p$ 不可約}，更準確地說，是絕對不可約的。$\bar E(x,y)$ 也是如此，因此模 $\mathfrak p$ 時同樣不會發生次數降低。

如果 $\bar E(x)$ 的係數是代數分數，還必須要求 $\mathfrak p$ 不同於整除公分母 $\Delta$ 的有限多個素理想；模其餘所有素理想時，可以用 $\Delta\cdot\bar E(x)$ 代替 $\bar E(x)$，從而滿足上述假設。定理得證。

\begin{flushleft}
埃朗根，1921 年 8 月。
\end{flushleft}
\begin{center}
（1921 年 8 月 28 日收到。）
\end{center}

\clearpage
\end{document}

